\tocsection{\textbf{10/7/26 Lecture 17:} Product Measures (cont.)}

Recall from last time we defined a measure on $(X \times, Y, \cal{A} \times \cal{B})$ with the property that
\[ (\mu \times \nu)(A \times B) = \mu(A) \times \nu(B) \]
and by Caratheodory, this is unique. Now, we need to define how to integrate with respect to this measure.\\

Let $E \subseteq X \times Y$, and, for fixed $x, y$, define
\[ E_{x} = \{y \in Y \col (x, y) \in E\} \]
and
\[ E^{y} = \{x \in X \col (x, y) \in E\} \]

Similarly, given a function $f : X \times Y \to \R$, for fixed $x, y$, define
\[ f_{x} : Y \to \R \tt{ by } f_{x}(y) = f(x, y) \]
and
\[ f^{y} : X \to \R \tt{ by } f^{y}(x) = f(x, y) \]

\begin{lemma}[\textbf{Proposition 17.1}]
    \vspace{-0.6cm}
    \begin{enumerate}
        \item If $E \in A \times B$, then $E_{x} \in \cal{B}$ for any $x \in X$ and $E^{y} \in \cal{A}$ for any $y \in Y$. 
        \item If $f$ is $\mu \times \nu$ measurable, then $f_{x}$ and $f^{y}$ are measurable with respect to $\mu$ and $\nu$.
    \end{enumerate}
\end{lemma}
\begin{proof}{Proof:}
    
    (1) Let $\cal{C}$ be the collection of subsets $E \subseteq X \times Y$ such that $E_{x} \in \cal{B}$ for all $x$ and $E^{y} \in \cal{A}$ for all $y$. We want to show that $\cal{C}$ contains measurable rectangles and is a $\sigma$-algebra.\\

    First, we see that if $E = A \times B$, then 
    \[ E_{x} = \left\{\begin{matrix} B & x \in A \\ \emptyset & x \notin A\end{matrix}\right. \]
    and
    \[ E^{y} = \left\{\begin{matrix} A & y \in B \\ \emptyset & y \notin B \end{matrix}\right. \]
    so all rectangles are in $\cal{C}$. Clearly, $\cal{C}$ is a $\sigma$-algebra as it is clear that it is closed under unions and complements (pass conditions to the restrictions). Thus, we see that $\cal{A} \times \cal{B} \in \cal{C}$, which was to be shown.\\

    (2) We see that for $U \subseteq \R$, so as
    \[ f_{x}^{-1}(U) = (f^{-1}(U))_{x}\]
    so if $f^{-1}(U) \in \cal{A} \times \cal{B}$, we see that $f_{x}^{-1}(U) \in \cal{B}$ by (1).

    $\qed$
\end{proof}

\begin{definition}[\textbf{Definition 17.2:} Monotone Classes]
    We say a set $\cal{M} \subseteq \cal{P}(X)$ is a \emph{monotone class} if
    \begin{enumerate}
        \item $A_1 \subseteq A_2 \subseteq \cdots,\quad A_n \in \cal{M} \implies \bigcup_{n} A_n \in \cal{M}$
        \item $A_1 \supseteq A_2 \supseteq \cdots,\quad A_n \in \cal{M} \implies \bigcap_{n} A_n \in \cal{M}$
    \end{enumerate} 
\end{definition}

Given a collection $\cal{C} \in \cal{P}(X)$, the monotone class generated by $\cal{C}$ is the smallest monotone class containing $\cal{C}$; that is, it is the intersection of all monotone classes containing $\cal{C}$. 

\begin{theorem}[\textbf{Theorem 17.3: }Monotone Class Theorem]
    Let $\cal{A} \subseteq \cal{P}(X)$ be an algebra. Let $\cal{M}$ be the monotone class generated by $\cal{A}$. Then, $M = \sigma(\cal{A})$. 
\end{theorem}
\begin{proof}{Proof:}
    Firstly, it is clear that $\cal{M} \subseteq \sigma(\cal{A})$ as $\sigma$-algebras are monotone classes. Now, we want to show that $\sigma(\cal{A}) \subseteq \cal{M}$, which is to show that $\cal{M}$ is a $\sigma$-algebra.\\
    
    Firstly, we need to show that if $A \in \cal{M}$, then $A^{c} \in \cal{M}$. Let $\cal{N}_1 = \{A \in \cal{M} \col A^{c} \in \cal{M}\}$. We see that $\cal{N}_1$ contains $\cal{A}$. Now, if $A_1 \subseteq A_2 \subseteq \cdots$ with $A_n \in \cal{N}_1$, since $\bigcup_{n} A_n \in \cal{M}$, and
    \[ \left(\bigcup_{n} A_n\right)^{c} = \left(\bigcap_{n} A_n^{c}\right) \]
    then as $A_1^{c} \supseteq A_2^{c} \supseteq \cdots$, we see that $\bigcap_{n} A_n^{c} = \left(\bigcup_{n} A_n\right)^{c} \in \cal{M}$. Thus, $\bigcup_{n} A_n \in \cal{N}_1$. By the same argument, we see that if $A_1 \supseteq A_2 \supseteq \cdots$, then $\bigcap_{n} A_n \in \cal{N}_1$, so $\cal{N}_1$ is a monotone class. Thus, $\cal{M} \subseteq \cal{N}_1$, so $\cal{M}$ is closed under complements.\\

    Now, we want to show that it is closed under intersection. Start by letting $\cal{N}_2 = \{B \in \cal{M} \col A \cap B \in \cal{M} \text{ for all } A \in \cal{A}\}$. We see that $\cal{N}_2$ contains $\cal{A}$. Now, suppose that $B_1 \subseteq B_2 \subseteq \cdots$ with $B_n \in \cal{N}_2$. Then, for any $A \in \cal{A}$, we see that as $\bigcup_{n} B_n \in \cal{M}$,
    \[ \bigcup_{n} B_n \cap A = \bigcup_{n} (B_n \cap A) \]
    and as $B_1 \cap A \subseteq B_2 \cap A \cdots$ and $B_n \cap A \in \cal{M}$, we see that $\bigcup_{n} B_n \cap A \in \cal{M}$, so $\bigcup_{n} B_n \in \cal{N}_2$. Since a similar proof for intersections follows, we see that $\cal{M} \subseteq \cal{N}_2$.\\

    For the next step, let $\cal{N}_3 = \{B \in \cal{M} \col A \cap B \in \cal{M} \tt{ for all } A \in \cal{M}\}$. Now, $\cal{N}_3$ contains $\cal{A}$ because of the last step. By the same argument as above, we see that $\cal{N}_3$ is a monotone class and so $\cal{M} \subseteq \cal{N}_3$.\\
    
    The last thing to check is countable intersections. Let $\{A_n\} \subseteq \cal{M}$. Now, since we know that $\bigcap_{i = 1}^{n} A_i \in \cal{M}$ for each $n$, and as
    \[ \bigcap_{i = 1}^{1} A_i \supseteq \bigcap_{i = 1}^{2} A_i \supseteq \cdots \]
    we see that
    \[ \bigcap_{n = 1}^{\infty} \bigcap_{i = 1}^{n} A_i = \bigcap_{i = 1}^{\infty} A_n \in \cal{M} \]
    Thus, $\cal{M}$ is a $\sigma$-algebra, which was to be proved.

    $\qed$
\end{proof}

\begin{lemma}[\textbf{Proposition 17.4}]
    Suppose that $\mu$ and $\nu$ are $\sigma$-finite. If $E \in \cal{A} \times \cal{B}$, then $x \mapsto \nu(E_{x})$ and $y \mapsto \mu(E^{y})$ are measurable on $X$ and $Y$ respectviely, and
    \[ (\mu \times \nu)(E) = \int_{X} \nu(E_{x})d\mu = \int_{Y} \mu(E^{y})dy \]
\end{lemma}
\begin{proof}{Proof:}
    First, assume that $\mu$ and $\nu$ are finite. Then, let $\cal{M}$ be the collection of the sets so that the conclusions of the theorem holds. We want to show that $\cal{M}$ contains finite unions of measurable rectangles and is a monotone class.\\

    First, let $E = A \times B$. Then, as
    \[ E_{x} = \left\{\begin{matrix} \emptyset & x \notin A \\ B & x \in A \end{matrix}\right. \]
    so $\nu(E_{x}) = \chi_{A}(x) \cdot \nu(B)$, which is measurable, and similarly for $E^{y}$. Similarly, we also see that
    \[ \int_{X} \nu(E_{x})dx = \mu(A)\nu(B) = (\mu \times \nu)(E) \]
    and similarly for the other integral. Thus, measurable rectangles are in $\cal{M}$.\\

    Now, let $E_1 \subseteq E_2 \subseteq \cdots$ are in $\cal{M}$ and let $E = \bigcup_{n} E_n$. Let $f_{n}(y) = \mu(E_{n}^{y})$. Each of these are measurable, and they increase pointwise to $f(y) = \mu(E^{y})$. Thus, $f$ is measurable as it is the limit of measurable functions. Then, by the monotone convergence theorem
    \begin{align*}
        \int_{Y} f(y) d\nu &= \lim_{n \to \infty} \int f_{n}(y)d\nu\\
        &= \lim_{n \to \infty} (\mu \times \nu)(E_n)\\
        &= (\mu \times \nu)(E)
    \end{align*}
    To see that $E_1 \supseteq E_2 \supseteq \cdots$ implies $E = \bigcap_{n} E_n$, it is the same proof but we use DCT using the fact ath $\mu \times \nu$ is $\sigma$-finite. 

    $\qed$
\end{proof}